% !TEX program = xelatex
% 伽马函数与贝塔函数
\documentclass[fontset=windows, chinese, 10pt, xcolor=dvipsnames]{beamer}
% Shared preamble for the chapter 11 standalone lectures.
% Each lecture sets \title and \subtitle after this file is read.

\usepackage{amsmath,amssymb,mathtools}
\usepackage{xeCJK}
\usepackage{zhnumber}
\usepackage{tikz}
\usepackage{tikz-3dplot}
\usetikzlibrary{arrows.meta,calc,positioning}

\usetheme{metropolis}
\usepackage{beamercolorthemeHistNanjing}

\newcommand{\red}[1]{\textcolor{red}{#1}}
\newcommand{\blue}[1]{\textcolor{blue}{#1}}
\newcommand{\green}[1]{\textcolor{green}{#1}}
\newcommand{\gray}[1]{\textcolor{gray}{#1}}
\newcommand{\orange}[1]{\textcolor{orange}{#1}}

\newcommand{\sssub}[1]{ {\mathrm{#1} }}  % script size subscript\\\
\newcommand{\funcd}[2]{ \frac{d{#1}}{d{#2}}} %function full derivative
\newcommand{\funcpd}[2]{ \frac{\partial{#1}}{\partial{#2}}} %function partial derivative
\newcommand{\fnald}[2] { \frac{\delta {#1}}{\delta {#2}}} % functional derivative

\renewcommand{\Re}{\operatorname{Re}}
\renewcommand{\Im}{\operatorname{Im}}
\def\ext{\sssub{ext}}
\def\x{\sssub{x}}
\def\xc{\sssub{xc}}
\def\H{\sssub{H}}
\def\Hx{\sssub{Hx}}
\def\Hxc{\sssub{Hxc}}
\def\c{\sssub{c}}
\def\s{\sssub{s}}

% \def\ext{\sssub{ext}}
% \def\ext{\sssub{ext}}
% \def\ext{\sssub{ext}}
% \def\ext{\sssub{ext}}
% \defsssub{ext}


% definitions of bold faced symbols
\def\brpppp{{\mathbf{r}^{\prime\prime\prime\prime}}}
\def\brppp{{\mathbf{r}^{\prime\prime\prime}}}
\def\brpp{{\mathbf{r}^{\prime\prime}}}
\def\brp{{\mathbf{r}^{\prime}}}
\def\bzp{{\mathbf{z}^{\prime}}}
\def\bxp{{\mathbf{x}^{\prime}}}
\def\tp{{{t}^{\prime}}}
\def\tpp{{{t}^{\prime\prime}}}
\def\tppp{{{t}^{\prime\prime\prime}}}

\def\tbr{{\tilde{\mathbf{r}}}}
\def\bk{{\mathbf{k}}}
\def\br{{\mathbf{r}}}
\def\bp{{\mathbf{p}}}
\def\bq{{\mathbf{q}}}
\def\bv{{\mathbf{v}}}
\def\bz{{\mathbf{z}}}
\def\mbf{{\mathbf{f}}}

\def\bx{{\mathbf{x}}}
\def\bR{{\mathbf{R}}}
\def\bM{{\mathbf{M}}}
\def\bP{{\mathbf{P}}}
\def\bT{{\mathbf{T}}}
\def\bK{{\mathbf{K}}}
\def\bA{{\mathbf{A}}}
\def\bB{{\mathbf{B}}}
\def\bD{{\mathbf{D}}}
\def\bE{{\mathbf{E}}}
\def\bF{{\mathbf{F}}}
\def\bG{{\mathbf{G}}}
\def\bH{{\mathbf{H}}}
\def\bI{{\mathbf{I}}}
\def\bJ{{\mathbf{J}}}
\def\bP{{\mathbf{P}}}
\def\bS{{\mathbf{S}}}
\def\bV{{\mathbf{V}}}
\def\bX{{\mathbf{X}}}
\def\bY{{\mathbf{Y}}}
\def\dd{{\mathop{}\!\mathrm{d}}}
\def\dr{{\!\!\dd\br\,}}
\def\dx{{\!\!\dd\bx\,}}
\def\dy{{\!\!\dd\by\,}}
\def\dk{{\!\!\dd\bk\,}}
\def\dG{{\!\!\dd\bG\,}}


\def\half{{\frac{1}{2}}}
\def\tr{{\mathrm{tr}}}


\def\Arg{{\mathrm{Arg}\;}}
\def\arg{{\mathrm{arg}\;}}
\def\Res{{\mathrm{Res}\;}}

% unit-vector basis (bold under XeLaTeX; \mathbf can look non-bold)
\newcommand{\ue}{\boldsymbol{e}}
\newcommand{\un}{\hat{\boldsymbol{n}}}


\author{罗凯}
\institute{南京理工大学物理学院}
\date{\zhdigits{\the\year},\zhganzhinian{\the\year}年秋}

\usebackgroundtemplate{%
\begin{tikzpicture}[remember picture,overlay]
  \node[opacity=0.05, at=(current page.center)] {
    \includegraphics[width=0.6\paperwidth]{njust-logo.jpeg}
  };
\end{tikzpicture}}


\title{特殊函数简介}
\subtitle{第一讲：伽马函数与贝塔函数}

\begin{document}

\begin{frame}
  \titlepage
\end{frame}

\begin{frame}{本讲提纲}
  阶乘只对正整数有定义。伽马函数（Gamma function）把它延拓到复数，贝塔函数（Beta function）是两个伽马函数的商。
  \begin{enumerate}
    \item 极限、积分、无穷乘积三种定义；
    \item 递推、反射公式和倍乘公式；
    \item 贝塔函数与伽马函数的关系；
    \item 麦克斯韦分布里的平均能量。
  \end{enumerate}
\end{frame}

\section{伽马函数}

\begin{frame}{欧拉的极限}
  \begin{equation*}
    \Gamma(z)
    =\lim_{n\to\infty}
      \frac{1\cdot 2\cdots n}{z(z+1)\cdots(z+n)}\,n^{z},
  \end{equation*}
  $z\neq 0,-1,-2,\ldots$。$z$ 可以是复数。
  \vspace{0.4em}
  把 $z$ 换成 $z+1$，分子多一个 $n$，分母多一个 $z+n+1$，
  \[
    \Gamma(z+1)
    =\lim_{n\to\infty}\frac{nz}{z+n+1}\,\Gamma(z)
    =z\,\Gamma(z).
  \]
\end{frame}

\begin{frame}{阶乘}
  \[
    \Gamma(1)
    =\lim_{n\to\infty}\frac{n!}{1\cdot 2\cdots n\cdot(n+1)}\,n
    =1.
  \]
  递推 $n-1$ 次，
  \begin{equation*}
    \boxed{\displaystyle
      \Gamma(n)=(n-1)!,
      \qquad n=1,2,3,\ldots.}
  \end{equation*}
  $\Gamma(2)=1$，$\Gamma(3)=2$，$\Gamma(4)=6$。
\end{frame}

\begin{frame}{欧拉积分}
  $\operatorname{Re} z>0$ 时，
  \begin{equation*}
    \boxed{\displaystyle
      \Gamma(z)=\int_0^{\infty} e^{-t} t^{z-1}\,\mathrm{d}t.}
  \end{equation*}
  分部积分立刻回到 $\Gamma(z+1)=z\Gamma(z)$。实部不大于零时，用这个递推向左延拓。
  \vspace{0.35em}
  常用的两个换元是
  \[
    \Gamma(z)=2\int_0^{\infty} e^{-t^2} t^{2z-1}\,\mathrm{d}t
    =\int_0^{1}\Bigl(\ln\frac{1}{t}\Bigr)^{z-1}\,\mathrm{d}t.
  \]
\end{frame}

\begin{frame}{$\Gamma(1/2)$}
  高斯积分 $\int_{-\infty}^{\infty}e^{-t^2}\,\mathrm{d}t=\sqrt{\pi}$。上面的换元在 $z=1/2$ 时只积正半轴，
  \[
    \Gamma\Bigl(\frac12\Bigr)
    =2\int_0^{\infty} e^{-t^2}\,\mathrm{d}t
    =\sqrt{\pi}.
  \]
  递推给出半整数，
  \[
    \Gamma\Bigl(\frac32\Bigr)=\frac12\sqrt{\pi},
    \qquad
    \Gamma\Bigl(\frac52\Bigr)=\frac34\sqrt{\pi}.
  \]
\end{frame}

\begin{frame}{魏尔斯特拉斯乘积}
  \[
    \frac{1}{\Gamma(z)}
    =z\,e^{\gamma z}
      \prod_{n=1}^{\infty}\Bigl(1+\frac{z}{n}\Bigr)e^{-z/n}.
  \]
  $\gamma=0.57721\ldots$ 是欧拉--马斯凯罗尼常数（Euler--Mascheroni constant）。
  \vspace{0.4em}
  乘积在 $z=0,-1,-2,\ldots$ 有一阶零点，所以 $\Gamma(z)$ 在这些点有一阶极点，其余地方解析。
\end{frame}

\begin{frame}{极点上的留数}
  $\Gamma(z+1)=z\Gamma(z)$，所以 $z=0$ 附近 $\Gamma(z)\sim 1/z$，留数是 $1$。
  \vspace{0.3em}
  用递推把 $\Gamma(-n+\varepsilon)$ 推到 $\Gamma(1+\varepsilon)$，分母是
  \[
    (-n+\varepsilon)(-n+1+\varepsilon)\cdots\varepsilon.
  \]
  $\varepsilon\to 0$ 时它变成 $(-1)^n n!\,\varepsilon$，因此
  \begin{equation*}
    \operatorname{Res}_{z=-n}\Gamma(z)=\frac{(-1)^n}{n!},
    \qquad n=0,1,2,\ldots.
  \end{equation*}
\end{frame}

\begin{frame}{反射公式}
  \begin{equation*}
    \boxed{\displaystyle
      \Gamma(z)\Gamma(1-z)=\frac{\pi}{\sin\pi z}.}
  \end{equation*}
  $z=1/2$ 时两边都是 $\pi$，和 $\Gamma(1/2)=\sqrt{\pi}$ 一致。
  \vspace{0.35em}
  它说明正半轴上的值决定负半轴：$\Gamma(1-z)$ 由 $\Gamma(z)$ 和正弦给出。极点正好落在正弦的零点上。
\end{frame}

\begin{frame}{倍乘公式}
  \begin{equation*}
    \boxed{\displaystyle
      \Gamma(1+z)\,\Gamma\Bigl(z+\frac12\Bigr)
      =2^{-2z}\sqrt{\pi}\,\Gamma(2z+1).}
  \end{equation*}
  等价的写法是
  \[
    \Gamma(z)\,\Gamma\Bigl(z+\frac12\Bigr)
    =2^{1-2z}\sqrt{\pi}\,\Gamma(2z).
  \]
  半整数的伽马函数都可以用它化成 $\sqrt{\pi}$ 乘一个有理数。
\end{frame}

\section{贝塔函数}

\begin{frame}{定义}
  $\operatorname{Re} p>0$，$\operatorname{Re} q>0$ 时，
  \begin{equation*}
    B(p,q)=\int_0^{1} t^{p-1}(1-t)^{q-1}\,\mathrm{d}t.
  \end{equation*}
  令 $t=\sin^2\theta$，
  \[
    B(p,q)
    =2\int_0^{\pi/2}
      \sin^{2p-1}\theta\,\cos^{2q-1}\theta\,\mathrm{d}\theta.
  \]
  $B(p,q)=B(q,p)$。
\end{frame}

\begin{frame}{化成伽马函数}
  把两个高斯型积分写成
  \[
    \Gamma(p)\Gamma(q)
    =4\int_0^{\infty}\!\!\int_0^{\infty}
      s^{2p-1}t^{2q-1}e^{-(s^2+t^2)}\,\mathrm{d}s\,\mathrm{d}t.
  \]
  极坐标 $s=r\cos\theta$、$t=r\sin\theta$，径向积出 $\Gamma(p+q)$，角度积出 $B(p,q)$。于是
  \begin{equation*}
    \boxed{\displaystyle
      B(p,q)=\frac{\Gamma(p)\Gamma(q)}{\Gamma(p+q)}.}
  \end{equation*}
\end{frame}

\begin{frame}{反射公式从这里出来}
  \[
    B(z,1-z)
    =\int_0^{1} t^{z-1}(1-t)^{-z}\,\mathrm{d}t.
  \]
  令 $x=t/(1-t)$，积分变成
  \[
    \int_0^{\infty}\frac{x^{z-1}}{1+x}\,\mathrm{d}x
    =\frac{\pi}{\sin\pi z}.
  \]
  左边又是 $\Gamma(z)\Gamma(1-z)$，这就是反射公式。倍乘公式同样由 $B(z,z)$ 作代换 $t=(1-\sqrt{\xi})/2$ 得到。
\end{frame}

\section{麦克斯韦分布}

\begin{frame}{归一化}
  理想气体的态密度正比于 $E^{1/2}$。记 $\beta=1/kT$，占据概率里的常数由
  \[
    1=C\int_0^{\infty} E^{1/2}e^{-\beta E}\,\mathrm{d}E
  \]
  定出。令 $u=\beta E$，
  \[
    \int_0^{\infty} E^{1/2}e^{-\beta E}\,\mathrm{d}E
    =\beta^{-3/2}\Gamma\Bigl(\frac32\Bigr)
    =\frac{\sqrt{\pi}}{2}\,\beta^{-3/2}.
  \]
  因此 $C=2\beta^{3/2}/\sqrt{\pi}$。
\end{frame}

\begin{frame}{平均能量}
  \[
    \langle E\rangle
    =C\int_0^{\infty} E^{3/2}e^{-\beta E}\,\mathrm{d}E
    =C\,\beta^{-5/2}\Gamma\Bigl(\frac52\Bigr).
  \]
  $\Gamma(5/2)=\frac12\cdot\frac32\sqrt{\pi}$，代入 $C$，
  \[
    \langle E\rangle
    =\frac{2\beta^{3/2}}{\sqrt{\pi}}
      \cdot\frac{3\sqrt{\pi}}{4}\,\beta^{-5/2}
    =\frac{3}{2\beta}
    =\frac32 kT.
  \]
  每个二次方自由度贡献 $\frac12 kT$，平动有三个。
\end{frame}

\begin{frame}{这一讲留下的公式}
  \[
    \Gamma(z+1)=z\Gamma(z),
    \qquad
    \Gamma(n)=(n-1)!,
    \qquad
    \Gamma\Bigl(\frac12\Bigr)=\sqrt{\pi}.
  \]
  \[
    \Gamma(z)\Gamma(1-z)=\frac{\pi}{\sin\pi z},
    \qquad
    B(p,q)=\frac{\Gamma(p)\Gamma(q)}{\Gamma(p+q)}.
  \]
  $z=-n$ 处的留数是 $(-1)^n/n!$。
\end{frame}

\end{document}
